\section{绪论}

\subsection{研究背景}

视觉语言模型的发展为多模态信息处理带来了很大推动。CLIP 等对比式模型通过海量图文数据的联合训练，首次展示了视觉和语言在统一嵌入空间中实现对齐的可能性；BLIP-2、LLaVA 和 Qwen-VL 等生成式模型在此基础上增加了语言解码能力，能够根据图像输入生成描述文本、回答问题和进行多轮对话。然而，这些模型在实际使用中的一个共同问题是：当输入图像出现模糊、低分辨率、噪声或遮挡时，模型输出的语义判断和文本生成质量会明显下降。

与此同时，脑电图作为一种成熟的脑机接口技术，能够以毫秒级的时间分辨率记录视觉刺激在大脑中引发的电生理响应。近年来有不少研究尝试从 EEG 信号中解码视觉信息——包括识别物体类别、检索语义相似的图像，以及将 EEG 表示映射到 CLIP 等预训练模型的嵌入空间。这让我们产生了一个想法：当机器视觉系统因为图像质量下降而看不清时，人脑的神经响应能否充当一种额外的信息源？

带着这个问题，我们围绕视觉退化条件下的 EEG 辅助视觉语言理解开展了本次课程设计。需要说明的是，我们的目标不是从脑电中直接解码语言，而是将 EEG 作为一种视觉语义增强信号来使用。在所有实验结论中，我们都通过与 shuffled EEG 和 random EEG 的对比来保证判断的客观性。

\subsection{课程设计目标与意义}

本次课程设计的目标是：在 Thought2Text 小样本 EEG-image 配对数据集上，搭建一个视觉退化条件下的 EEG 辅助视觉语言理解系统，并通过系统的对照实验来评估真实 EEG 能否在六种退化场景下改善语义识别和图像描述生成的性能。

从学习实践的角度看，这个课题让我们有机会把课堂上接触到的深度学习、Transformer、多模态融合、大模型微调等知识点串联起来，在一个真实数据集上完整经历从问题定义、数据处理、模型搭建、训练调参到结果分析的整个流程。特别是在单 GPU 环境下完成大模型适配、处理 EEG 这种非标准模态数据，带来了很多工程层面的收获。

从方法验证的角度看，我们设置了四组对照（vision-only、real EEG、shuffled EEG、random EEG），贯穿所有实验评估。这一设计确保任何观察到的性能提升都可以追溯到 EEG 和图像之间的真实配对关系，而不是模型容量增加或输入通道增多带来的附带效应。

\subsection{研究内容与技术路线}

本次课程设计主要完成了以下工作：

\begin{enumerate}
  \item 搭建了包含六种退化条件和四种 EEG 对照模式的完整实验框架，采用 image-level split 防止训练-测试信号泄漏；
  \item 设计并实现了 A2 时间-频谱-空间 EEG 语义融合模型，从 EEG 的时间、频谱和空间通道三个维度并行提取特征，通过门控残差融合对退化图像进行语义补充；
  \item 实现了 VTF token 级融合方法，将 EEG token 通过交叉注意力注入到 CLIP ViT 的视觉 patch token 中，生成增强视觉表示；
  \item 构建了基于 Qwen2-VL 的生成式图像描述模块，使用 LoRA 进行参数高效微调，对比了五种生成方案，并通过多候选重排序提升输出质量。
\end{enumerate}

我们的技术路线采用了从简到繁的思路：先用 CLIP 的预训练语义空间作为跨模态桥梁，降低对 EEG 数据量的依赖；再在受约束的分类任务上验证 EEG 语义增强的有效性；确认效果后再推进到 token 级融合和自由文本生成。每完成一步都有明确的中间检查结果，出现问题也方便定位。

\subsection{报告结构安排}

本报告共分十章。第 1 章绪论；第 2 章相关技术基础；第 3 章数据集与任务；第 4 章系统总体设计；第 5 章 A2 EEG 语义融合模型；第 6 章 token 级 EEG 视觉增强；第 7 章生成式视觉语言模型设计；第 8 章实验结果与分析；第 9 章工程实现与代码；第 10 章总结与反思。
